%%%%%%%%%%%%%%%%%%%%%%%%%%%%%%%%%%%%%%%%%%%%%%%%%%%%%%%%%%%%
% 3.10 SYMPLECTIC GEOMETRY
% The Composition of Advanced Mathematics
%%%%%%%%%%%%%%%%%%%%%%%%%%%%%%%%%%%%%%%%%%%%%%%%%%%%%%%%%%%%

\section{Symplectic Geometry}
\label{sec:symplectic-geometry}

\prerequisite{
Linear Algebra, Multivariable Calculus, Differential Forms,
Differential Geometry from \cref{sec:differential-geometry},
and Complex Geometry from \cref{sec:complex-geometry}.
}

\learningobjectives{
By the end of this section, the reader should be able to:
\begin{itemize}
    \item explain the basic purpose of symplectic geometry;
    \item define a symplectic form and a symplectic manifold;
    \item verify that a differential \(2\)-form is closed and nondegenerate;
    \item explain why symplectic manifolds must have even dimension;
    \item work with the standard symplectic form on \(\R^{2n}\);
    \item recognize canonical coordinates \((q_i,p_i)\);
    \item explain symplectic vector spaces and symplectic bases;
    \item define symplectomorphisms;
    \item state Darboux's theorem and explain its geometric significance;
    \item define Hamiltonian vector fields;
    \item derive Hamilton's equations from the symplectic form;
    \item define and compute Poisson brackets;
    \item explain conserved quantities in Hamiltonian systems;
    \item describe cotangent bundles and their canonical symplectic structures;
    \item distinguish isotropic, coisotropic, and Lagrangian subspaces;
    \item recognize Lagrangian submanifolds;
    \item explain symplectic area and volume;
    \item describe the relationship between symplectic and K\"ahler geometry;
    \item recognize moment maps and symplectic reduction at an introductory level;
    \item explain major applications of symplectic geometry.
\end{itemize}
}

%%%%%%%%%%%%%%%%%%%%%%%%%%%%%%%%%%%%%%%%%%%%%%%%%%%%%%%%%%%%
\subsection{What Is Symplectic Geometry?}
\label{subsec:what-is-symplectic-geometry}
%%%%%%%%%%%%%%%%%%%%%%%%%%%%%%%%%%%%%%%%%%%%%%%%%%%%%%%%%%%%

Symplectic geometry studies spaces equipped with a special differential
\(2\)-form that measures oriented two-dimensional behavior.

Its origins lie in classical mechanics, where the state of a mechanical
system is described by both position and momentum.

For one particle moving along a line, the state is represented by

\[
(q,p),
\]

where

\[
q=\text{position},
\qquad
p=\text{momentum}.
\]

For \(n\) degrees of freedom, the natural coordinates are

\[
\boxed{
(q_1,\ldots,q_n,p_1,\ldots,p_n).
}
\]

The resulting space has dimension \(2n\).

\begin{conceptbox}
Riemannian geometry is fundamentally concerned with lengths and angles.

Symplectic geometry is fundamentally concerned with oriented areas and the
geometry of Hamiltonian motion.
\end{conceptbox}

%%%%%%%%%%%%%%%%%%%%%%%%%%%%%%%%%%%%%%%%%%%%%%%%%%%%%%%%%%%%
\subsection{Differential Two-Forms}
\label{subsec:symplectic-two-forms}
%%%%%%%%%%%%%%%%%%%%%%%%%%%%%%%%%%%%%%%%%%%%%%%%%%%%%%%%%%%%

At each point \(p\) of a smooth manifold \(M\), a differential \(2\)-form
assigns an alternating bilinear map

\[
\omega_p:T_pM\times T_pM\to\R.
\]

Thus

\[
\boxed{
\omega_p(u,v)=-\omega_p(v,u).
}
\]

In particular,

\[
\boxed{
\omega_p(v,v)=0.
}
\]

The lowercase Greek letter

\[
\omega
\]

is read ``omega.''

In symplectic geometry, \(\omega\) is usually called the
\emph{symplectic form}.

%%%%%%%%%%%%%%%%%%%%%%%%%%%%%%%%%%%%%%%%%%%%%%%%%%%%%%%%%%%%
\subsection{Symplectic Forms}
\label{subsec:symplectic-forms}
%%%%%%%%%%%%%%%%%%%%%%%%%%%%%%%%%%%%%%%%%%%%%%%%%%%%%%%%%%%%

\begin{definitionbox}{Symplectic Form}
A differential \(2\)-form

\[
\omega\in\Omega^2(M)
\]

is a \emph{symplectic form} if it satisfies two conditions:

\[
\boxed{d\omega=0}
\]

and

\[
\boxed{\omega\text{ is nondegenerate}.}
\]
\end{definitionbox}

The first condition means that \(\omega\) is \emph{closed}.

The second means that if

\[
\omega_p(v,w)=0
\]

for every

\[
w\in T_pM,
\]

then

\[
\boxed{v=0.}
\]

%%%%%%%%%%%%%%%%%%%%%%%%%%%%%%%%%%%%%%%%%%%%%%%%%%%%%%%%%%%%
\subsection{Symplectic Manifolds}
\label{subsec:symplectic-manifolds}
%%%%%%%%%%%%%%%%%%%%%%%%%%%%%%%%%%%%%%%%%%%%%%%%%%%%%%%%%%%%

\begin{definitionbox}{Symplectic Manifold}
A \emph{symplectic manifold} is a pair

\[
\boxed{(M,\omega)}
\]

where \(M\) is a smooth manifold and \(\omega\) is a symplectic form on
\(M\).
\end{definitionbox}

Thus symplectic geometry studies manifolds carrying a closed,
nondegenerate \(2\)-form.

%%%%%%%%%%%%%%%%%%%%%%%%%%%%%%%%%%%%%%%%%%%%%%%%%%%%%%%%%%%%
\subsection{The Standard Symplectic Form}
\label{subsec:standard-symplectic-form}
%%%%%%%%%%%%%%%%%%%%%%%%%%%%%%%%%%%%%%%%%%%%%%%%%%%%%%%%%%%%

The fundamental example is

\[
\R^{2n}
\]

with coordinates

\[
(q_1,\ldots,q_n,p_1,\ldots,p_n).
\]

The standard symplectic form is

\[
\boxed{
\omega_0
=
\sum_{j=1}^{n}
dq_j\wedge dp_j.
}
\]

For

\[
n=1,
\]

this becomes

\[
\boxed{
\omega_0=dq\wedge dp.
}
\]

Thus the basic symplectic plane is the position-momentum plane.

%%%%%%%%%%%%%%%%%%%%%%%%%%%%%%%%%%%%%%%%%%%%%%%%%%%%%%%%%%%%
\subsection{Position and Momentum Coordinates}
\label{subsec:position-momentum-coordinates}
%%%%%%%%%%%%%%%%%%%%%%%%%%%%%%%%%%%%%%%%%%%%%%%%%%%%%%%%%%%%

The variables

\[
q_j
\]

are traditionally called \emph{position coordinates}, while

\[
p_j
\]

are called \emph{momentum coordinates}.

The pair

\[
(q_j,p_j)
\]

is called a pair of \emph{canonical coordinates}.

The notation reflects the origin of symplectic geometry in Hamiltonian
mechanics.

\begin{centereddiagram}
\begin{tikzpicture}[scale=1.05]
\draw[-{Latex[length=3mm]}] (-3,0)--(3.3,0)
node[right] {$q$};
\draw[-{Latex[length=3mm]}] (0,-2.3)--(0,2.5)
node[above] {$p$};

\draw[fill=gray!15,thick]
(.5,.4) rectangle (2.2,1.6);

\draw[<->] (.5,.15)--(2.2,.15);
\node at (1.35,-.1) {$\Delta q$};

\draw[<->] (2.45,.4)--(2.45,1.6);
\node[right] at (2.45,1) {$\Delta p$};

\node at (1.35,1) {$\Delta q\,\Delta p$};
\end{tikzpicture}
\end{centereddiagram}

The form

\[
dq\wedge dp
\]

measures oriented area in this plane.

%%%%%%%%%%%%%%%%%%%%%%%%%%%%%%%%%%%%%%%%%%%%%%%%%%%%%%%%%%%%
\subsection{Closedness of the Standard Form}
\label{subsec:standard-form-closed}
%%%%%%%%%%%%%%%%%%%%%%%%%%%%%%%%%%%%%%%%%%%%%%%%%%%%%%%%%%%%

Consider

\[
\omega_0
=
\sum_{j=1}^{n}
dq_j\wedge dp_j.
\]

Since

\[
d(dq_j)=0
\]

and

\[
d(dp_j)=0,
\]

we obtain

\[
d(dq_j\wedge dp_j)=0.
\]

Therefore

\[
\boxed{
d\omega_0=0.
}
\]

Hence the standard symplectic form is closed.

%%%%%%%%%%%%%%%%%%%%%%%%%%%%%%%%%%%%%%%%%%%%%%%%%%%%%%%%%%%%
\subsection{Nondegeneracy}
\label{subsec:symplectic-nondegeneracy}
%%%%%%%%%%%%%%%%%%%%%%%%%%%%%%%%%%%%%%%%%%%%%%%%%%%%%%%%%%%%

Nondegeneracy means that the symplectic form detects every nonzero tangent
vector.

If

\[
v\neq0,
\]

then there must exist some vector \(w\) such that

\[
\boxed{
\omega(v,w)\neq0.
}
\]

Equivalently, the map

\[
\boxed{
v\longmapsto\omega(v,\cdot)
}
\]

from the tangent space to its dual is an isomorphism.

%%%%%%%%%%%%%%%%%%%%%%%%%%%%%%%%%%%%%%%%%%%%%%%%%%%%%%%%%%%%
\subsection{Matrix Representation}
\label{subsec:symplectic-matrix-representation}
%%%%%%%%%%%%%%%%%%%%%%%%%%%%%%%%%%%%%%%%%%%%%%%%%%%%%%%%%%%%

On \(\R^{2n}\), the standard symplectic form can be represented by the
matrix

\[
\boxed{
J=
\begin{pmatrix}
0&I_n\\
-I_n&0
\end{pmatrix}.
}
\]

For vectors \(u,v\in\R^{2n}\),

\[
\boxed{
\omega_0(u,v)=u^T Jv.
}
\]

The matrix \(J\) satisfies

\[
\boxed{
J^T=-J
}
\]

and

\[
\boxed{
J^2=-I_{2n}.
}
\]

\begin{notebox}
The matrix \(J\) resembles the standard complex structure from
\cref{sec:complex-geometry}. This is one of the first indications of the
close relationship between symplectic and complex geometry.
\end{notebox}

%%%%%%%%%%%%%%%%%%%%%%%%%%%%%%%%%%%%%%%%%%%%%%%%%%%%%%%%%%%%
\subsection{Symplectic Vector Spaces}
\label{subsec:symplectic-vector-spaces}
%%%%%%%%%%%%%%%%%%%%%%%%%%%%%%%%%%%%%%%%%%%%%%%%%%%%%%%%%%%%

\begin{definitionbox}{Symplectic Vector Space}
A \emph{symplectic vector space} is a real vector space \(V\) equipped with
a nondegenerate alternating bilinear form

\[
\boxed{
\omega:V\times V\to\R.
}
\]
\end{definitionbox}

This is the linear-algebraic model for the tangent spaces of a symplectic
manifold.

At every point \(p\in M\),

\[
\boxed{
(T_pM,\omega_p)
}
\]

is a symplectic vector space.

%%%%%%%%%%%%%%%%%%%%%%%%%%%%%%%%%%%%%%%%%%%%%%%%%%%%%%%%%%%%
\subsection{Why Symplectic Spaces Have Even Dimension}
\label{subsec:symplectic-even-dimension}
%%%%%%%%%%%%%%%%%%%%%%%%%%%%%%%%%%%%%%%%%%%%%%%%%%%%%%%%%%%%

A nondegenerate skew-symmetric matrix must have even size.

Therefore a symplectic vector space must have even dimension.

Consequently,

\[
\boxed{
\dim M=2n
}
\]

for every symplectic manifold.

There are no symplectic manifolds of odd dimension.

\begin{importantbox}
Symplectic geometry is intrinsically even-dimensional.

The basic dimensions are

\[
2,\ 4,\ 6,\ 8,\ldots
\]
\end{importantbox}

%%%%%%%%%%%%%%%%%%%%%%%%%%%%%%%%%%%%%%%%%%%%%%%%%%%%%%%%%%%%
\subsection{Symplectic Bases}
\label{subsec:symplectic-bases}
%%%%%%%%%%%%%%%%%%%%%%%%%%%%%%%%%%%%%%%%%%%%%%%%%%%%%%%%%%%%

A symplectic vector space of dimension \(2n\) admits a basis

\[
e_1,\ldots,e_n,f_1,\ldots,f_n
\]

such that

\[
\boxed{
\omega(e_i,e_j)=0,
}
\]

\[
\boxed{
\omega(f_i,f_j)=0,
}
\]

and

\[
\boxed{
\omega(e_i,f_j)=\delta_{ij}.
}
\]

Here

\[
\delta_{ij}
\]

is the Kronecker delta.

Such a basis is called a \emph{symplectic basis}.

%%%%%%%%%%%%%%%%%%%%%%%%%%%%%%%%%%%%%%%%%%%%%%%%%%%%%%%%%%%%
\subsection{Symplectic Linear Transformations}
\label{subsec:symplectic-linear-transformations}
%%%%%%%%%%%%%%%%%%%%%%%%%%%%%%%%%%%%%%%%%%%%%%%%%%%%%%%%%%%%

A linear transformation

\[
A:\R^{2n}\to\R^{2n}
\]

is symplectic when it preserves the standard symplectic form.

In matrix form,

\[
\boxed{
A^TJA=J.
}
\]

The collection of such matrices forms the \emph{symplectic group}

\[
\boxed{
\operatorname{Sp}(2n,\R).
}
\]

%%%%%%%%%%%%%%%%%%%%%%%%%%%%%%%%%%%%%%%%%%%%%%%%%%%%%%%%%%%%
\subsection{Symplectomorphisms}
\label{subsec:symplectomorphisms}
%%%%%%%%%%%%%%%%%%%%%%%%%%%%%%%%%%%%%%%%%%%%%%%%%%%%%%%%%%%%

The nonlinear analogue of a symplectic linear transformation is a
symplectomorphism.

\begin{definitionbox}{Symplectomorphism}
Let

\[
(M,\omega_M)
\]

and

\[
(N,\omega_N)
\]

be symplectic manifolds.

A diffeomorphism

\[
\phi:M\to N
\]

is a \emph{symplectomorphism} if

\[
\boxed{
\phi^*\omega_N=\omega_M.
}
\]
\end{definitionbox}

Thus a symplectomorphism preserves the symplectic structure.

Two symplectic manifolds related by a symplectomorphism are regarded as
equivalent in symplectic geometry.

%%%%%%%%%%%%%%%%%%%%%%%%%%%%%%%%%%%%%%%%%%%%%%%%%%%%%%%%%%%%
\subsection{Darboux's Theorem}
\label{subsec:darboux-theorem}
%%%%%%%%%%%%%%%%%%%%%%%%%%%%%%%%%%%%%%%%%%%%%%%%%%%%%%%%%%%%

One of the most striking facts in symplectic geometry is that every
symplectic manifold looks locally the same.

\begin{theorem}[Darboux's Theorem]
Let \((M,\omega)\) be a symplectic manifold of dimension \(2n\). Around every
point there exist local coordinates

\[
(q_1,\ldots,q_n,p_1,\ldots,p_n)
\]

such that

\[
\boxed{
\omega
=
\sum_{j=1}^{n}
dq_j\wedge dp_j.
}
\]
\end{theorem}

These are called \emph{Darboux coordinates}.

%%%%%%%%%%%%%%%%%%%%%%%%%%%%%%%%%%%%%%%%%%%%%%%%%%%%%%%%%%%%
\subsection{Why Darboux's Theorem Is Remarkable}
\label{subsec:darboux-significance}
%%%%%%%%%%%%%%%%%%%%%%%%%%%%%%%%%%%%%%%%%%%%%%%%%%%%%%%%%%%%

In Riemannian geometry, curvature provides local invariants.

Two Riemannian manifolds may have very different local geometry.

Darboux's theorem says that symplectic geometry has no analogous local
invariants for the symplectic form alone.

Locally,

\[
\boxed{
\text{every symplectic manifold looks like }
(\R^{2n},\omega_0).
}
\]

Therefore much of the richness of symplectic geometry is global rather than
purely local.

%%%%%%%%%%%%%%%%%%%%%%%%%%%%%%%%%%%%%%%%%%%%%%%%%%%%%%%%%%%%
\subsection{Phase Space}
\label{subsec:phase-space}
%%%%%%%%%%%%%%%%%%%%%%%%%%%%%%%%%%%%%%%%%%%%%%%%%%%%%%%%%%%%

In classical mechanics, the collection of all positions and momenta of a
system is called its \emph{phase space}.

For \(n\) degrees of freedom,

\[
\boxed{
\text{phase space}
=
\{(q_1,\ldots,q_n,p_1,\ldots,p_n)\}.
}
\]

A point in phase space represents the complete instantaneous mechanical
state of the system.

The evolution of the system traces a curve through phase space.

%%%%%%%%%%%%%%%%%%%%%%%%%%%%%%%%%%%%%%%%%%%%%%%%%%%%%%%%%%%%
\subsection{Hamiltonian Functions}
\label{subsec:hamiltonian-functions}
%%%%%%%%%%%%%%%%%%%%%%%%%%%%%%%%%%%%%%%%%%%%%%%%%%%%%%%%%%%%

A smooth function

\[
\boxed{
H:M\to\R
}
\]

on a symplectic manifold is called a \emph{Hamiltonian}.

In mechanics,

\[
H
\]

usually represents the total energy of the system.

For example,

\[
\boxed{
H(q,p)
=
\frac{p^2}{2m}
+
V(q)
}
\]

describes a particle of mass \(m\) moving in a potential \(V(q)\).

%%%%%%%%%%%%%%%%%%%%%%%%%%%%%%%%%%%%%%%%%%%%%%%%%%%%%%%%%%%%
\subsection{Hamiltonian Vector Fields}
\label{subsec:hamiltonian-vector-fields}
%%%%%%%%%%%%%%%%%%%%%%%%%%%%%%%%%%%%%%%%%%%%%%%%%%%%%%%%%%%%

Because \(\omega\) is nondegenerate, every smooth function \(H\) determines
a unique vector field.

\begin{definitionbox}{Hamiltonian Vector Field}
The \emph{Hamiltonian vector field} \(X_H\) associated with a smooth
function \(H\) is defined by

\[
\boxed{
\iota_{X_H}\omega=-dH.
}
\]
\end{definitionbox}

Here

\[
\iota_{X_H}
\]

denotes contraction, also called the interior product, with the vector field
\(X_H\).

%%%%%%%%%%%%%%%%%%%%%%%%%%%%%%%%%%%%%%%%%%%%%%%%%%%%%%%%%%%%
\subsection{Deriving Hamilton's Equations}
\label{subsec:deriving-hamilton-equations}
%%%%%%%%%%%%%%%%%%%%%%%%%%%%%%%%%%%%%%%%%%%%%%%%%%%%%%%%%%%%

On

\[
\R^{2n},
\]

write

\[
X_H
=
\sum_{j=1}^{n}
\left(
a_j\frac{\partial}{\partial q_j}
+
b_j\frac{\partial}{\partial p_j}
\right).
\]

Using

\[
\omega
=
\sum_{j=1}^{n}
dq_j\wedge dp_j,
\]

we obtain

\[
\iota_{X_H}\omega
=
\sum_{j=1}^{n}
\left(
a_j\,dp_j-b_j\,dq_j
\right).
\]

Meanwhile,

\[
-dH
=
-\sum_{j=1}^{n}
\left(
\frac{\partial H}{\partial q_j}dq_j
+
\frac{\partial H}{\partial p_j}dp_j
\right).
\]

Matching coefficients gives

\[
a_j
=
-\frac{\partial H}{\partial p_j},
\qquad
b_j
=
\frac{\partial H}{\partial q_j}
\]

under the convention

\[
\iota_{X_H}\omega=-dH.
\]

If instead one uses the common mechanics convention

\[
\omega
=
\sum_{j=1}^{n}
dp_j\wedge dq_j,
\]

the familiar signs reverse.

To maintain the standard Hamilton equations below, we adopt the equivalent
convention

\[
\boxed{
\iota_{X_H}\omega=dH
}
\]

for

\[
\omega=\sum_j dq_j\wedge dp_j.
\]

Then

\[
\boxed{
X_H
=
\sum_{j=1}^{n}
\left(
\frac{\partial H}{\partial p_j}
\frac{\partial}{\partial q_j}
-
\frac{\partial H}{\partial q_j}
\frac{\partial}{\partial p_j}
\right).
}
\]

\begin{notebox}
Sign conventions for Hamiltonian vector fields differ among textbooks.
The essential point is to choose a convention and use it consistently.
In the remainder of this section we use

\[
\iota_{X_H}\omega=dH
\]

with

\[
\omega=\sum_j dq_j\wedge dp_j.
\]
\end{notebox}

%%%%%%%%%%%%%%%%%%%%%%%%%%%%%%%%%%%%%%%%%%%%%%%%%%%%%%%%%%%%
\subsection{Hamilton's Equations}
\label{subsec:hamiltons-equations}
%%%%%%%%%%%%%%%%%%%%%%%%%%%%%%%%%%%%%%%%%%%%%%%%%%%%%%%%%%%%

The integral curves of \(X_H\) satisfy

\[
\boxed{
\dot q_j
=
\frac{\partial H}{\partial p_j}
}
\]

and

\[
\boxed{
\dot p_j
=
-
\frac{\partial H}{\partial q_j}.
}
\]

These are \emph{Hamilton's equations}.

Together,

\[
\boxed{
\begin{aligned}
\dot q_j&=\frac{\partial H}{\partial p_j},\\
\dot p_j&=-\frac{\partial H}{\partial q_j}.
\end{aligned}
}
\]

They describe the time evolution of a Hamiltonian system.

%%%%%%%%%%%%%%%%%%%%%%%%%%%%%%%%%%%%%%%%%%%%%%%%%%%%%%%%%%%%
\subsection{Example: Particle in a Potential}
\label{subsec:hamiltonian-particle-potential}
%%%%%%%%%%%%%%%%%%%%%%%%%%%%%%%%%%%%%%%%%%%%%%%%%%%%%%%%%%%%

Consider

\[
H(q,p)
=
\frac{p^2}{2m}
+
V(q).
\]

Then

\[
\frac{\partial H}{\partial p}
=
\frac{p}{m}
\]

and

\[
\frac{\partial H}{\partial q}
=
V'(q).
\]

Hamilton's equations give

\[
\boxed{
\dot q=\frac{p}{m}
}
\]

and

\[
\boxed{
\dot p=-V'(q).
}
\]

Since

\[
p=m\dot q,
\]

we obtain

\[
m\ddot q=-V'(q),
\]

which is Newton's second law for a conservative force.

%%%%%%%%%%%%%%%%%%%%%%%%%%%%%%%%%%%%%%%%%%%%%%%%%%%%%%%%%%%%
\subsection{Example: Harmonic Oscillator}
\label{subsec:symplectic-harmonic-oscillator}
%%%%%%%%%%%%%%%%%%%%%%%%%%%%%%%%%%%%%%%%%%%%%%%%%%%%%%%%%%%%

For a harmonic oscillator,

\[
V(q)=\frac12kq^2.
\]

The Hamiltonian is

\[
\boxed{
H(q,p)
=
\frac{p^2}{2m}
+
\frac12kq^2.
}
\]

Hamilton's equations are

\[
\boxed{
\dot q=\frac{p}{m}
}
\]

and

\[
\boxed{
\dot p=-kq.
}
\]

Combining them gives

\[
\boxed{
m\ddot q+kq=0.
}
\]

The trajectories in phase space lie on energy curves

\[
\frac{p^2}{2m}
+
\frac12kq^2
=
E,
\]

which are ellipses.

\begin{centereddiagram}
\begin{tikzpicture}[scale=1]
\draw[-{Latex[length=3mm]}] (-3,0)--(3.2,0)
node[right] {$q$};
\draw[-{Latex[length=3mm]}] (0,-2.2)--(0,2.4)
node[above] {$p$};

\draw[thick] (0,0) ellipse (2.3 and 1.4);
\draw[thick] (0,0) ellipse (1.5 and .9);
\draw[thick] (0,0) ellipse (.75 and .45);

\draw[-{Latex[length=2.5mm]}]
(2.3,0) arc[start angle=0,end angle=30,x radius=2.3,y radius=1.4];

\node at (2.1,1.55) {$H=E$};
\end{tikzpicture}
\end{centereddiagram}

%%%%%%%%%%%%%%%%%%%%%%%%%%%%%%%%%%%%%%%%%%%%%%%%%%%%%%%%%%%%
\subsection{Hamiltonian Flow}
\label{subsec:hamiltonian-flow}
%%%%%%%%%%%%%%%%%%%%%%%%%%%%%%%%%%%%%%%%%%%%%%%%%%%%%%%%%%%%

The vector field \(X_H\) generates a family of transformations

\[
\phi_t:M\to M.
\]

This family is called the \emph{Hamiltonian flow} of \(H\).

The parameter

\[
t
\]

usually represents time.

Hamiltonian flow preserves the symplectic form:

\[
\boxed{
\phi_t^*\omega=\omega.
}
\]

Therefore Hamiltonian evolution is symplectic.

%%%%%%%%%%%%%%%%%%%%%%%%%%%%%%%%%%%%%%%%%%%%%%%%%%%%%%%%%%%%
\subsection{Conservation of the Hamiltonian}
\label{subsec:hamiltonian-conservation}
%%%%%%%%%%%%%%%%%%%%%%%%%%%%%%%%%%%%%%%%%%%%%%%%%%%%%%%%%%%%

Along a trajectory generated by \(H\),

\[
\frac{dH}{dt}
=
dH(X_H).
\]

Using

\[
dH=\iota_{X_H}\omega,
\]

we obtain

\[
dH(X_H)
=
\omega(X_H,X_H).
\]

Since \(\omega\) is alternating,

\[
\omega(X_H,X_H)=0.
\]

Therefore

\[
\boxed{
\frac{dH}{dt}=0.
}
\]

Thus the Hamiltonian is conserved along its own flow.

%%%%%%%%%%%%%%%%%%%%%%%%%%%%%%%%%%%%%%%%%%%%%%%%%%%%%%%%%%%%
\subsection{Poisson Brackets}
\label{subsec:poisson-brackets}
%%%%%%%%%%%%%%%%%%%%%%%%%%%%%%%%%%%%%%%%%%%%%%%%%%%%%%%%%%%%

Two smooth functions can be combined using the symplectic structure.

\begin{definitionbox}{Poisson Bracket}
For smooth functions \(f,g\), define the \emph{Poisson bracket} by

\[
\boxed{
\{f,g\}
=
\sum_{j=1}^{n}
\left(
\frac{\partial f}{\partial q_j}
\frac{\partial g}{\partial p_j}
-
\frac{\partial f}{\partial p_j}
\frac{\partial g}{\partial q_j}
\right).
}
\]
\end{definitionbox}

The notation

\[
\{f,g\}
\]

is read ``the Poisson bracket of \(f\) and \(g\).''

With the convention used here,

\[
\boxed{
\{f,g\}
=
\omega(X_f,X_g).
}
\]

%%%%%%%%%%%%%%%%%%%%%%%%%%%%%%%%%%%%%%%%%%%%%%%%%%%%%%%%%%%%
\subsection{Canonical Poisson Relations}
\label{subsec:canonical-poisson-relations}
%%%%%%%%%%%%%%%%%%%%%%%%%%%%%%%%%%%%%%%%%%%%%%%%%%%%%%%%%%%%

For canonical coordinates,

\[
\boxed{
\{q_i,q_j\}=0,
}
\]

\[
\boxed{
\{p_i,p_j\}=0,
}
\]

and

\[
\boxed{
\{q_i,p_j\}=\delta_{ij}.
}
\]

These relations encode the fundamental interaction between position and
momentum variables.

%%%%%%%%%%%%%%%%%%%%%%%%%%%%%%%%%%%%%%%%%%%%%%%%%%%%%%%%%%%%
\subsection{Properties of the Poisson Bracket}
\label{subsec:poisson-bracket-properties}
%%%%%%%%%%%%%%%%%%%%%%%%%%%%%%%%%%%%%%%%%%%%%%%%%%%%%%%%%%%%

The Poisson bracket is bilinear:

\[
\boxed{
\{af+bg,h\}
=
a\{f,h\}+b\{g,h\}.
}
\]

It is antisymmetric:

\[
\boxed{
\{f,g\}=-\{g,f\}.
}
\]

It satisfies the product rule:

\[
\boxed{
\{f,gh\}
=
\{f,g\}h+g\{f,h\}.
}
\]

It also satisfies the Jacobi identity:

\[
\boxed{
\{f,\{g,h\}\}
+
\{g,\{h,f\}\}
+
\{h,\{f,g\}\}
=
0.
}
\]

Therefore smooth functions on a symplectic manifold form a Lie algebra under
the Poisson bracket.

%%%%%%%%%%%%%%%%%%%%%%%%%%%%%%%%%%%%%%%%%%%%%%%%%%%%%%%%%%%%
\subsection{Time Evolution and Poisson Brackets}
\label{subsec:time-evolution-poisson}
%%%%%%%%%%%%%%%%%%%%%%%%%%%%%%%%%%%%%%%%%%%%%%%%%%%%%%%%%%%%

Let \(f\) be an observable along a Hamiltonian trajectory.

Then

\[
\boxed{
\frac{df}{dt}
=
\{f,H\}
+
\frac{\partial f}{\partial t}.
}
\]

If \(f\) has no explicit time dependence,

\[
\boxed{
\frac{df}{dt}
=
\{f,H\}.
}
\]

Therefore \(f\) is conserved when

\[
\boxed{
\{f,H\}=0.
}
\]

%%%%%%%%%%%%%%%%%%%%%%%%%%%%%%%%%%%%%%%%%%%%%%%%%%%%%%%%%%%%
\subsection{Cotangent Bundles}
\label{subsec:symplectic-cotangent-bundles}
%%%%%%%%%%%%%%%%%%%%%%%%%%%%%%%%%%%%%%%%%%%%%%%%%%%%%%%%%%%%

Let \(Q\) be a smooth manifold.

Its cotangent bundle

\[
\boxed{
T^*Q
}
\]

consists of covectors attached to points of \(Q\).

In mechanics,

\[
Q
\]

is interpreted as configuration space, while

\[
T^*Q
\]

is phase space.

If

\[
\dim Q=n,
\]

then

\[
\boxed{
\dim T^*Q=2n.
}
\]

%%%%%%%%%%%%%%%%%%%%%%%%%%%%%%%%%%%%%%%%%%%%%%%%%%%%%%%%%%%%
\subsection{The Canonical One-Form}
\label{subsec:canonical-one-form}
%%%%%%%%%%%%%%%%%%%%%%%%%%%%%%%%%%%%%%%%%%%%%%%%%%%%%%%%%%%%

Every cotangent bundle carries a natural \(1\)-form.

In local coordinates

\[
(q_1,\ldots,q_n,p_1,\ldots,p_n),
\]

define

\[
\boxed{
\theta
=
\sum_{j=1}^{n}
p_j\,dq_j.
}
\]

The lowercase Greek letter

\[
\theta
\]

is read ``theta.''

This form is called the \emph{canonical one-form}, \emph{tautological
one-form}, or \emph{Liouville one-form}.

%%%%%%%%%%%%%%%%%%%%%%%%%%%%%%%%%%%%%%%%%%%%%%%%%%%%%%%%%%%%
\subsection{Canonical Symplectic Form on \(T^*Q\)}
\label{subsec:canonical-cotangent-symplectic-form}
%%%%%%%%%%%%%%%%%%%%%%%%%%%%%%%%%%%%%%%%%%%%%%%%%%%%%%%%%%%%

Using the convention of this section, define

\[
\boxed{
\omega=-d\theta.
}
\]

Since

\[
\theta
=
\sum_j p_j\,dq_j,
\]

we have

\[
d\theta
=
\sum_j dp_j\wedge dq_j.
\]

Therefore

\[
\boxed{
\omega
=
\sum_j dq_j\wedge dp_j.
}
\]

Thus every cotangent bundle possesses a canonical symplectic structure.

\begin{importantbox}
A manifold \(Q\) does not need to be symplectic for its cotangent bundle
\(T^*Q\) to be symplectic.

The symplectic structure on \(T^*Q\) is canonical.
\end{importantbox}

%%%%%%%%%%%%%%%%%%%%%%%%%%%%%%%%%%%%%%%%%%%%%%%%%%%%%%%%%%%%
\subsection{Symplectic Orthogonal Complements}
\label{subsec:symplectic-orthogonal}
%%%%%%%%%%%%%%%%%%%%%%%%%%%%%%%%%%%%%%%%%%%%%%%%%%%%%%%%%%%%

Let \(W\subseteq V\) be a subspace of a symplectic vector space.

Define its symplectic orthogonal complement by

\[
\boxed{
W^\omega
=
\{
v\in V:
\omega(v,w)=0
\text{ for all }w\in W
\}.
}
\]

This resembles the ordinary orthogonal complement from inner-product
geometry, but it is defined using the symplectic form rather than a metric.

%%%%%%%%%%%%%%%%%%%%%%%%%%%%%%%%%%%%%%%%%%%%%%%%%%%%%%%%%%%%
\subsection{Isotropic Subspaces}
\label{subsec:isotropic-subspaces}
%%%%%%%%%%%%%%%%%%%%%%%%%%%%%%%%%%%%%%%%%%%%%%%%%%%%%%%%%%%%

\begin{definitionbox}{Isotropic Subspace}
A subspace \(W\subseteq V\) is \emph{isotropic} if

\[
\boxed{
W\subseteq W^\omega.
}
\]
\end{definitionbox}

Equivalently,

\[
\boxed{
\omega|_W=0.
}
\]

Thus every pair of vectors in \(W\) has zero symplectic pairing.

%%%%%%%%%%%%%%%%%%%%%%%%%%%%%%%%%%%%%%%%%%%%%%%%%%%%%%%%%%%%
\subsection{Coisotropic Subspaces}
\label{subsec:coisotropic-subspaces}
%%%%%%%%%%%%%%%%%%%%%%%%%%%%%%%%%%%%%%%%%%%%%%%%%%%%%%%%%%%%

\begin{definitionbox}{Coisotropic Subspace}
A subspace \(W\subseteq V\) is \emph{coisotropic} if

\[
\boxed{
W^\omega\subseteq W.
}
\]
\end{definitionbox}

Coisotropic subspaces play an important role in constraints and symplectic
reduction.

%%%%%%%%%%%%%%%%%%%%%%%%%%%%%%%%%%%%%%%%%%%%%%%%%%%%%%%%%%%%
\subsection{Lagrangian Subspaces}
\label{subsec:lagrangian-subspaces}
%%%%%%%%%%%%%%%%%%%%%%%%%%%%%%%%%%%%%%%%%%%%%%%%%%%%%%%%%%%%

\begin{definitionbox}{Lagrangian Subspace}
A subspace \(L\) of a \(2n\)-dimensional symplectic vector space is
\emph{Lagrangian} if

\[
\boxed{
L=L^\omega.
}
\]
\end{definitionbox}

Equivalently,

\[
\boxed{
\omega|_L=0
}
\]

and

\[
\boxed{
\dim L=n.
}
\]

Thus Lagrangian subspaces are maximal isotropic subspaces.

%%%%%%%%%%%%%%%%%%%%%%%%%%%%%%%%%%%%%%%%%%%%%%%%%%%%%%%%%%%%
\subsection{Example of a Lagrangian Subspace}
\label{subsec:lagrangian-subspace-example}
%%%%%%%%%%%%%%%%%%%%%%%%%%%%%%%%%%%%%%%%%%%%%%%%%%%%%%%%%%%%

Consider

\[
\R^{2n}
\]

with coordinates

\[
(q_1,\ldots,q_n,p_1,\ldots,p_n).
\]

The subspace

\[
\boxed{
L=
\{(q_1,\ldots,q_n,0,\ldots,0)\}
}
\]

is Lagrangian.

Its dimension is \(n\), and on \(L\),

\[
dp_j=0.
\]

Therefore

\[
\omega_0|_L=0.
\]

%%%%%%%%%%%%%%%%%%%%%%%%%%%%%%%%%%%%%%%%%%%%%%%%%%%%%%%%%%%%
\subsection{Lagrangian Submanifolds}
\label{subsec:lagrangian-submanifolds}
%%%%%%%%%%%%%%%%%%%%%%%%%%%%%%%%%%%%%%%%%%%%%%%%%%%%%%%%%%%%

\begin{definitionbox}{Lagrangian Submanifold}
Let \((M,\omega)\) have dimension \(2n\).

An \(n\)-dimensional submanifold

\[
L\subseteq M
\]

is \emph{Lagrangian} if

\[
\boxed{
\omega|_L=0.
}
\]
\end{definitionbox}

At every point \(p\in L\),

\[
T_pL
\]

is a Lagrangian subspace of

\[
T_pM.
\]

%%%%%%%%%%%%%%%%%%%%%%%%%%%%%%%%%%%%%%%%%%%%%%%%%%%%%%%%%%%%
\subsection{Graphs of Exact One-Forms}
\label{subsec:graphs-exact-one-forms}
%%%%%%%%%%%%%%%%%%%%%%%%%%%%%%%%%%%%%%%%%%%%%%%%%%%%%%%%%%%%

Let

\[
f:Q\to\R
\]

be smooth.

Its differential

\[
df
\]

defines a section of \(T^*Q\).

The graph

\[
\boxed{
L_f
=
\{(q,df_q):q\in Q\}
}
\]

is a Lagrangian submanifold of \(T^*Q\).

This gives a large and important family of Lagrangian submanifolds.

%%%%%%%%%%%%%%%%%%%%%%%%%%%%%%%%%%%%%%%%%%%%%%%%%%%%%%%%%%%%
\subsection{Symplectic Area}
\label{subsec:symplectic-area}
%%%%%%%%%%%%%%%%%%%%%%%%%%%%%%%%%%%%%%%%%%%%%%%%%%%%%%%%%%%%

On a two-dimensional symplectic manifold,

\[
\omega
\]

acts as an area form.

If \(S\) is an oriented region, its symplectic area is

\[
\boxed{
\int_S\omega.
}
\]

For the standard plane,

\[
\omega=dq\wedge dp,
\]

so

\[
\boxed{
\int_S dq\wedge dp
}
\]

is the ordinary signed area in phase space.

%%%%%%%%%%%%%%%%%%%%%%%%%%%%%%%%%%%%%%%%%%%%%%%%%%%%%%%%%%%%
\subsection{Symplectic Volume}
\label{subsec:symplectic-volume}
%%%%%%%%%%%%%%%%%%%%%%%%%%%%%%%%%%%%%%%%%%%%%%%%%%%%%%%%%%%%

On a \(2n\)-dimensional symplectic manifold,

\[
\omega^n
=
\underbrace{
\omega\wedge\cdots\wedge\omega
}_{n\text{ times}}
\]

is a nonzero top-degree form.

The natural volume form is

\[
\boxed{
\frac{\omega^n}{n!}.
}
\]

Thus a symplectic form automatically determines an orientation and a volume
measure.

%%%%%%%%%%%%%%%%%%%%%%%%%%%%%%%%%%%%%%%%%%%%%%%%%%%%%%%%%%%%
\subsection{Symplectic Manifolds Are Orientable}
\label{subsec:symplectic-orientability}
%%%%%%%%%%%%%%%%%%%%%%%%%%%%%%%%%%%%%%%%%%%%%%%%%%%%%%%%%%%%

Because

\[
\omega^n
\]

is nowhere zero on a \(2n\)-dimensional symplectic manifold, it determines
an orientation.

Therefore

\[
\boxed{
\text{every symplectic manifold is orientable}.
}
\]

This gives an immediate topological restriction on manifolds that can admit
symplectic structures.

%%%%%%%%%%%%%%%%%%%%%%%%%%%%%%%%%%%%%%%%%%%%%%%%%%%%%%%%%%%%
\subsection{Liouville's Theorem}
\label{subsec:liouville-theorem-symplectic}
%%%%%%%%%%%%%%%%%%%%%%%%%%%%%%%%%%%%%%%%%%%%%%%%%%%%%%%%%%%%

Hamiltonian flows preserve the symplectic form.

Consequently, they also preserve symplectic volume.

\begin{theorem}[Liouville's Theorem]
Hamiltonian flow preserves the phase-space volume form

\[
\boxed{
\frac{\omega^n}{n!}.
}
\]
\end{theorem}

Thus Hamiltonian evolution cannot compress an entire region of phase space
into a smaller phase-space volume.

%%%%%%%%%%%%%%%%%%%%%%%%%%%%%%%%%%%%%%%%%%%%%%%%%%%%%%%%%%%%
\subsection{Area Preservation in Dimension Two}
\label{subsec:area-preservation-symplectic}
%%%%%%%%%%%%%%%%%%%%%%%%%%%%%%%%%%%%%%%%%%%%%%%%%%%%%%%%%%%%

When

\[
n=1,
\]

symplectic volume is simply area.

Therefore a symplectomorphism of a symplectic surface preserves oriented
area.

Schematically,

\[
\boxed{
\text{symplectic map}
\Longrightarrow
\text{area preserving}
}
\]

in dimension two.

In higher dimensions, volume preservation alone is much weaker than
symplectic preservation.

%%%%%%%%%%%%%%%%%%%%%%%%%%%%%%%%%%%%%%%%%%%%%%%%%%%%%%%%%%%%
\subsection{Gromov's Non-Squeezing Theorem}
\label{subsec:gromov-nonsqueezing}
%%%%%%%%%%%%%%%%%%%%%%%%%%%%%%%%%%%%%%%%%%%%%%%%%%%%%%%%%%%%

Symplectic geometry has rigidity phenomena that cannot be detected merely by
volume.

One of the most famous examples is Gromov's non-squeezing theorem.

Very roughly, a symplectic ball of radius \(R\) cannot be symplectically
embedded into a symplectic cylinder of radius \(r\) when

\[
\boxed{
r<R.
}
\]

Even if the cylinder has enough total volume, the symplectic embedding may
still be impossible.

This phenomenon is sometimes called the \emph{symplectic camel} principle.

%%%%%%%%%%%%%%%%%%%%%%%%%%%%%%%%%%%%%%%%%%%%%%%%%%%%%%%%%%%%
\subsection{Symplectic Capacity}
\label{subsec:symplectic-capacity}
%%%%%%%%%%%%%%%%%%%%%%%%%%%%%%%%%%%%%%%%%%%%%%%%%%%%%%%%%%%%

The non-squeezing theorem motivates quantities called
\emph{symplectic capacities}.

A symplectic capacity assigns a number

\[
c(M,\omega)
\]

that measures symplectic size in a way compatible with symplectic
embeddings.

Unlike ordinary volume, a symplectic capacity detects the special
two-dimensional rigidity inherent in symplectic geometry.

%%%%%%%%%%%%%%%%%%%%%%%%%%%%%%%%%%%%%%%%%%%%%%%%%%%%%%%%%%%%
\subsection{Hamiltonian Group Actions}
\label{subsec:hamiltonian-group-actions}
%%%%%%%%%%%%%%%%%%%%%%%%%%%%%%%%%%%%%%%%%%%%%%%%%%%%%%%%%%%%

Suppose a Lie group \(G\) acts on a symplectic manifold

\[
(M,\omega)
\]

by symplectomorphisms.

Under additional conditions, the action may be described by a
\emph{moment map}.

Moment maps connect

\[
\boxed{
\text{symmetry}
+
\text{conservation laws}
+
\text{symplectic geometry}.
}
\]

%%%%%%%%%%%%%%%%%%%%%%%%%%%%%%%%%%%%%%%%%%%%%%%%%%%%%%%%%%%%
\subsection{Moment Maps}
\label{subsec:moment-maps}
%%%%%%%%%%%%%%%%%%%%%%%%%%%%%%%%%%%%%%%%%%%%%%%%%%%%%%%%%%%%

Let

\[
\mathfrak g
\]

denote the Lie algebra of \(G\), and let

\[
\mathfrak g^*
\]

be its dual.

A moment map has the form

\[
\boxed{
\mu:M\to\mathfrak g^*.
}
\]

The lowercase Greek letter

\[
\mu
\]

is read ``mu.''

For each

\[
\xi\in\mathfrak g,
\]

the corresponding infinitesimal group action is related to the Hamiltonian
function

\[
\langle\mu,\xi\rangle.
\]

%%%%%%%%%%%%%%%%%%%%%%%%%%%%%%%%%%%%%%%%%%%%%%%%%%%%%%%%%%%%
\subsection{Symplectic Reduction}
\label{subsec:symplectic-reduction}
%%%%%%%%%%%%%%%%%%%%%%%%%%%%%%%%%%%%%%%%%%%%%%%%%%%%%%%%%%%%

Symmetry can be used to construct smaller symplectic spaces.

Under suitable regularity and freeness assumptions, one forms

\[
\boxed{
M_{\mathrm{red}}
=
\mu^{-1}(c)/G_c,
}
\]

where \(c\) is a chosen value of the moment map and \(G_c\) is the
appropriate symmetry group preserving that level.

This procedure is called \emph{symplectic reduction} or
\emph{Marsden--Weinstein reduction}.

Conceptually,

\[
\boxed{
\text{constraints}
+
\text{symmetry quotient}
\longrightarrow
\text{reduced phase space}.
}
\]

%%%%%%%%%%%%%%%%%%%%%%%%%%%%%%%%%%%%%%%%%%%%%%%%%%%%%%%%%%%%
\subsection{Symplectic Geometry and Noether's Principle}
\label{subsec:symplectic-noether}
%%%%%%%%%%%%%%%%%%%%%%%%%%%%%%%%%%%%%%%%%%%%%%%%%%%%%%%%%%%%

Continuous symmetries of Hamiltonian systems are associated with conserved
quantities.

Examples include:

\[
\boxed{
\text{translation symmetry}
\longleftrightarrow
\text{momentum conservation},
}
\]

and

\[
\boxed{
\text{rotational symmetry}
\longleftrightarrow
\text{angular momentum conservation}.
}
\]

Moment maps provide the geometric language for these conservation laws.

%%%%%%%%%%%%%%%%%%%%%%%%%%%%%%%%%%%%%%%%%%%%%%%%%%%%%%%%%%%%
\subsection{Compatible Almost Complex Structures}
\label{subsec:symplectic-compatible-complex}
%%%%%%%%%%%%%%%%%%%%%%%%%%%%%%%%%%%%%%%%%%%%%%%%%%%%%%%%%%%%

Let

\[
(M,\omega)
\]

be symplectic.

An almost complex structure

\[
J:TM\to TM
\]

is compatible with \(\omega\) when

\[
\boxed{
\omega(JX,JY)=\omega(X,Y)
}
\]

and

\[
\boxed{
\omega(X,JX)>0
}
\]

for every nonzero \(X\).

Then

\[
\boxed{
g(X,Y)=\omega(X,JY)
}
\]

defines a Riemannian metric.

%%%%%%%%%%%%%%%%%%%%%%%%%%%%%%%%%%%%%%%%%%%%%%%%%%%%%%%%%%%%
\subsection{The Symplectic--Complex--Metric Relationship}
\label{subsec:symplectic-complex-metric}
%%%%%%%%%%%%%%%%%%%%%%%%%%%%%%%%%%%%%%%%%%%%%%%%%%%%%%%%%%%%

A compatible pair

\[
(\omega,J)
\]

produces a metric

\[
g.
\]

These structures satisfy

\[
\boxed{
g(X,Y)=\omega(X,JY).
}
\]

Thus the three objects

\[
\boxed{
g,\qquad J,\qquad\omega
}
\]

are tightly connected.

In K\"ahler geometry, \(J\) is additionally integrable.

%%%%%%%%%%%%%%%%%%%%%%%%%%%%%%%%%%%%%%%%%%%%%%%%%%%%%%%%%%%%
\subsection{K\"ahler Geometry Revisited}
\label{subsec:symplectic-kahler}
%%%%%%%%%%%%%%%%%%%%%%%%%%%%%%%%%%%%%%%%%%%%%%%%%%%%%%%%%%%%

A K\"ahler manifold has

\[
\boxed{
(M,g,J,\omega)
}
\]

with compatible

\[
\text{Riemannian metric }g,
\]

\[
\text{complex structure }J,
\]

and

\[
\text{symplectic form }\omega.
\]

Therefore

\[
\boxed{
\text{K\"ahler}
\Longrightarrow
\text{symplectic}.
}
\]

The converse is false.

A symplectic manifold need not admit an integrable compatible complex
structure.

%%%%%%%%%%%%%%%%%%%%%%%%%%%%%%%%%%%%%%%%%%%%%%%%%%%%%%%%%%%%
\subsection{Exact Symplectic Forms}
\label{subsec:exact-symplectic-forms}
%%%%%%%%%%%%%%%%%%%%%%%%%%%%%%%%%%%%%%%%%%%%%%%%%%%%%%%%%%%%

A symplectic form is called \emph{exact} if there exists a \(1\)-form
\(\lambda\) such that

\[
\boxed{
\omega=d\lambda.
}
\]

The standard form on \(\R^{2n}\) is exact.

For example,

\[
\lambda
=
\sum_{j=1}^{n}
q_j\,dp_j
\]

satisfies

\[
d\lambda
=
\sum_{j=1}^{n}
dq_j\wedge dp_j
=
\omega_0.
\]

%%%%%%%%%%%%%%%%%%%%%%%%%%%%%%%%%%%%%%%%%%%%%%%%%%%%%%%%%%%%
\subsection{Compact Symplectic Manifolds Are Not Exact}
\label{subsec:compact-symplectic-not-exact}
%%%%%%%%%%%%%%%%%%%%%%%%%%%%%%%%%%%%%%%%%%%%%%%%%%%%%%%%%%%%

Suppose \(M\) is compact without boundary and

\[
\omega=d\lambda.
\]

Then

\[
\omega^n
=
d(\lambda\wedge\omega^{n-1}),
\]

because

\[
d\omega=0.
\]

By Stokes' theorem,

\[
\int_M\omega^n=0.
\]

But \(\omega^n\) is a volume form and has nonzero integral with the
symplectic orientation.

This is impossible.

Therefore

\[
\boxed{
\text{a compact symplectic manifold without boundary cannot have an exact
symplectic form}.
}
\]

%%%%%%%%%%%%%%%%%%%%%%%%%%%%%%%%%%%%%%%%%%%%%%%%%%%%%%%%%%%%
\subsection{Contact Geometry Preview}
\label{subsec:contact-geometry-preview}
%%%%%%%%%%%%%%%%%%%%%%%%%%%%%%%%%%%%%%%%%%%%%%%%%%%%%%%%%%%%

Symplectic geometry naturally interacts with odd-dimensional geometry.

The odd-dimensional analogue is \emph{contact geometry}.

A contact manifold has dimension

\[
\boxed{2n+1}
\]

and is equipped locally with a \(1\)-form \(\alpha\) satisfying

\[
\boxed{
\alpha\wedge(d\alpha)^n\neq0.
}
\]

Thus

\[
\boxed{
\text{symplectic geometry}
\leftrightarrow
\text{even dimensions},
}
\]

while

\[
\boxed{
\text{contact geometry}
\leftrightarrow
\text{odd dimensions}.
}
\]

%%%%%%%%%%%%%%%%%%%%%%%%%%%%%%%%%%%%%%%%%%%%%%%%%%%%%%%%%%%%
\subsection{Worked Example: Standard Symplectic Pairing}
\label{subsec:worked-symplectic-pairing}
%%%%%%%%%%%%%%%%%%%%%%%%%%%%%%%%%%%%%%%%%%%%%%%%%%%%%%%%%%%%

\begin{workedexamplebox}{Compute a Symplectic Pairing}

On \(\R^2\), let

\[
\omega=dq\wedge dp.
\]

Consider

\[
u=
\begin{pmatrix}
2\\1
\end{pmatrix},
\qquad
v=
\begin{pmatrix}
3\\4
\end{pmatrix}.
\]

For vectors

\[
u=(u_q,u_p),
\qquad
v=(v_q,v_p),
\]

we have

\[
\omega(u,v)
=
u_qv_p-u_pv_q.
\]

Therefore

\[
\omega(u,v)
=
2(4)-1(3)
=
5.
\]

\begin{answerbox}
\[
\boxed{\omega(u,v)=5}
\]
\end{answerbox}
\end{workedexamplebox}

%%%%%%%%%%%%%%%%%%%%%%%%%%%%%%%%%%%%%%%%%%%%%%%%%%%%%%%%%%%%
\subsection{Worked Example: Testing Nondegeneracy}
\label{subsec:worked-symplectic-nondegeneracy}
%%%%%%%%%%%%%%%%%%%%%%%%%%%%%%%%%%%%%%%%%%%%%%%%%%%%%%%%%%%%

\begin{workedexamplebox}{Nondegeneracy on \(\R^2\)}

Let

\[
\omega=dq\wedge dp.
\]

Suppose

\[
v=(a,b)
\]

satisfies

\[
\omega(v,w)=0
\]

for every

\[
w=(c,d).
\]

Then

\[
ad-bc=0
\]

for all \(c,d\).

Choose

\[
(c,d)=(0,1).
\]

Then

\[
a=0.
\]

Choose

\[
(c,d)=(1,0).
\]

Then

\[
-b=0,
\]

so

\[
b=0.
\]

Therefore

\[
v=0.
\]

\begin{answerbox}
\[
\boxed{
dq\wedge dp
\text{ is nondegenerate on }\R^2.
}
\]
\end{answerbox}
\end{workedexamplebox}

%%%%%%%%%%%%%%%%%%%%%%%%%%%%%%%%%%%%%%%%%%%%%%%%%%%%%%%%%%%%
\subsection{Worked Example: Hamilton's Equations}
\label{subsec:worked-hamilton-equations}
%%%%%%%%%%%%%%%%%%%%%%%%%%%%%%%%%%%%%%%%%%%%%%%%%%%%%%%%%%%%

\begin{workedexamplebox}{Free Particle}

Consider

\[
H(q,p)=\frac{p^2}{2m}.
\]

Then

\[
\frac{\partial H}{\partial p}
=
\frac{p}{m}
\]

and

\[
\frac{\partial H}{\partial q}=0.
\]

Hamilton's equations give

\[
\dot q=\frac{p}{m},
\qquad
\dot p=0.
\]

Thus momentum is constant:

\[
p(t)=p_0.
\]

Therefore

\[
q(t)
=
q_0+\frac{p_0}{m}t.
\]

\begin{answerbox}
\[
\boxed{
p(t)=p_0,
\qquad
q(t)=q_0+\frac{p_0}{m}t
}
\]
\end{answerbox}
\end{workedexamplebox}

%%%%%%%%%%%%%%%%%%%%%%%%%%%%%%%%%%%%%%%%%%%%%%%%%%%%%%%%%%%%
\subsection{Worked Example: Poisson Bracket}
\label{subsec:worked-poisson-bracket}
%%%%%%%%%%%%%%%%%%%%%%%%%%%%%%%%%%%%%%%%%%%%%%%%%%%%%%%%%%%%

\begin{workedexamplebox}{Compute a Poisson Bracket}

Let

\[
f(q,p)=q^2
\]

and

\[
g(q,p)=p^2.
\]

Then

\[
\frac{\partial f}{\partial q}=2q,
\qquad
\frac{\partial f}{\partial p}=0,
\]

and

\[
\frac{\partial g}{\partial q}=0,
\qquad
\frac{\partial g}{\partial p}=2p.
\]

Therefore

\[
\{f,g\}
=
(2q)(2p)-(0)(0).
\]

Hence

\[
\begin{answerbox}
\[
\boxed{
\{q^2,p^2\}=4qp
}
\]
\end{answerbox}
\end{workedexamplebox}

%%%%%%%%%%%%%%%%%%%%%%%%%%%%%%%%%%%%%%%%%%%%%%%%%%%%%%%%%%%%
\subsection{Worked Example: A Lagrangian Submanifold}
\label{subsec:worked-lagrangian}
%%%%%%%%%%%%%%%%%%%%%%%%%%%%%%%%%%%%%%%%%%%%%%%%%%%%%%%%%%%%

\begin{workedexamplebox}{The \(q\)-Axis}

Consider

\[
\R^2
\]

with

\[
\omega=dq\wedge dp.
\]

Let

\[
L=\{(q,0):q\in\R\}.
\]

The submanifold \(L\) has dimension

\[
1=\frac12\dim\R^2.
\]

Along \(L\),

\[
p=0
\]

and tangent vectors have no \(p\)-component.

Therefore

\[
\omega|_L=0.
\]

Hence \(L\) is Lagrangian.

\begin{answerbox}
\[
\boxed{
L=\{(q,0)\}
\text{ is Lagrangian.}
}
\]
\end{answerbox}
\end{workedexamplebox}

%%%%%%%%%%%%%%%%%%%%%%%%%%%%%%%%%%%%%%%%%%%%%%%%%%%%%%%%%%%%
\subsection{Common Errors}
\label{subsec:symplectic-common-errors}
%%%%%%%%%%%%%%%%%%%%%%%%%%%%%%%%%%%%%%%%%%%%%%%%%%%%%%%%%%%%

\subsubsection{Thinking a Symplectic Form Is a Metric}

A metric satisfies

\[
g(v,v)>0
\]

for nonzero \(v\).

A symplectic form satisfies

\[
\boxed{
\omega(v,v)=0
}
\]

for every \(v\).

Thus a symplectic form does not directly measure vector lengths.

%%%%%%%%%%%%%%%%%%%%%%%%%%%%%%%%%%%%%%%%%%%%%%%%%%%%%%%%%%%%
\subsubsection{Forgetting Closedness}

Nondegeneracy alone is not enough.

A symplectic form must satisfy both

\[
\boxed{
d\omega=0
}
\]

and

\[
\boxed{
\omega\text{ nondegenerate}.
}
\]

%%%%%%%%%%%%%%%%%%%%%%%%%%%%%%%%%%%%%%%%%%%%%%%%%%%%%%%%%%%%
\subsubsection{Confusing Closed and Exact}

If

\[
\omega=d\lambda,
\]

then automatically

\[
d\omega=0.
\]

Thus

\[
\boxed{
\text{exact}
\Longrightarrow
\text{closed}.
}
\]

The converse is not generally true.

%%%%%%%%%%%%%%%%%%%%%%%%%%%%%%%%%%%%%%%%%%%%%%%%%%%%%%%%%%%%
\subsubsection{Trying to Put a Symplectic Form on an Odd-Dimensional Manifold}

A symplectic manifold must have even dimension.

Therefore spaces of dimension

\[
1,\ 3,\ 5,\ldots
\]

cannot themselves be symplectic.

%%%%%%%%%%%%%%%%%%%%%%%%%%%%%%%%%%%%%%%%%%%%%%%%%%%%%%%%%%%%
\subsubsection{Confusing Volume Preservation with Symplectic Preservation}

Every symplectomorphism preserves symplectic volume.

However,

\[
\boxed{
\text{volume preserving}
\not\Longrightarrow
\text{symplectic}
}
\]

in dimensions greater than two.

%%%%%%%%%%%%%%%%%%%%%%%%%%%%%%%%%%%%%%%%%%%%%%%%%%%%%%%%%%%%
\subsubsection{Confusing \(q\) and \(p\)}

In Hamiltonian mechanics,

\[
q_j
\]

usually denotes position and

\[
p_j
\]

usually denotes momentum.

Their roles are paired but not interchangeable without changing the
symplectic conventions.

%%%%%%%%%%%%%%%%%%%%%%%%%%%%%%%%%%%%%%%%%%%%%%%%%%%%%%%%%%%%
\subsubsection{Forgetting Hamiltonian Sign Conventions}

Different authors define Hamiltonian vector fields using either

\[
\iota_{X_H}\omega=dH
\]

or

\[
\iota_{X_H}\omega=-dH.
\]

The resulting formulas differ by signs.

Always check the convention being used.

%%%%%%%%%%%%%%%%%%%%%%%%%%%%%%%%%%%%%%%%%%%%%%%%%%%%%%%%%%%%
\subsubsection{Confusing Isotropic and Lagrangian}

Every Lagrangian subspace is isotropic.

However, an isotropic subspace need not be Lagrangian.

A Lagrangian subspace of a \(2n\)-dimensional symplectic vector space must
have dimension exactly

\[
\boxed{n.}
\]

%%%%%%%%%%%%%%%%%%%%%%%%%%%%%%%%%%%%%%%%%%%%%%%%%%%%%%%%%%%%
\subsubsection{Assuming Every Symplectic Manifold Is K\"ahler}

Every K\"ahler manifold is symplectic.

The converse is false.

A symplectic manifold does not automatically possess an integrable complex
structure compatible with its symplectic form.

%%%%%%%%%%%%%%%%%%%%%%%%%%%%%%%%%%%%%%%%%%%%%%%%%%%%%%%%%%%%
\subsubsection{Thinking Darboux's Theorem Makes Symplectic Geometry Trivial}

Darboux's theorem says that symplectic manifolds are locally standard.

It does not say they are globally equivalent.

Global topology, Hamiltonian dynamics, Lagrangian submanifolds, and
symplectic embeddings can be highly nontrivial.

%%%%%%%%%%%%%%%%%%%%%%%%%%%%%%%%%%%%%%%%%%%%%%%%%%%%%%%%%%%%
\subsection{Applications and Connections}
\label{subsec:symplectic-applications}
%%%%%%%%%%%%%%%%%%%%%%%%%%%%%%%%%%%%%%%%%%%%%%%%%%%%%%%%%%%%

\begin{applicationbox}
\textbf{Classical Mechanics.}
Symplectic manifolds provide the natural geometric setting for Hamiltonian
mechanics. Position and momentum form canonical coordinate pairs, and
Hamilton's equations describe motion as symplectic flow.

\medskip

\textbf{Celestial Mechanics.}
Planetary and orbital systems can be studied using Hamiltonian dynamics,
conservation laws, perturbation theory, and symplectic methods.

\medskip

\textbf{Dynamical Systems.}
Hamiltonian systems form a major class of dynamical systems. Periodic
orbits, invariant sets, stability, and integrability are deeply connected
with symplectic structure.

\medskip

\textbf{Differential Geometry.}
Symplectic forms are differential forms satisfying strong nondegeneracy
conditions, connecting the subject directly with smooth manifolds and
differential topology.

\medskip

\textbf{Complex Geometry.}
K\"ahler manifolds combine symplectic forms with compatible complex
structures and Riemannian metrics.

\medskip

\textbf{Algebraic Geometry.}
Complex projective varieties provide important symplectic manifolds, and
symplectic methods have become important in enumerative and algebraic
geometry.

\medskip

\textbf{Topology.}
Symplectic structures impose global topological restrictions, while
symplectic invariants distinguish spaces that may appear similar from a
purely smooth viewpoint.

\medskip

\textbf{Mathematical Physics.}
Phase spaces, conservation laws, geometric quantization, gauge theory, and
string theory all involve symplectic structures.

\medskip

\textbf{Optimization and Control.}
Hamiltonian systems and related geometric structures appear in optimal
control, variational problems, and mechanics-based optimization.

\medskip

\textbf{Numerical Mathematics.}
Symplectic integrators are numerical methods designed to preserve the
geometric structure of Hamiltonian systems over long time intervals.
\end{applicationbox}

%%%%%%%%%%%%%%%%%%%%%%%%%%%%%%%%%%%%%%%%%%%%%%%%%%%%%%%%%%%%
\subsection{Exercises}
\label{subsec:symplectic-exercises}
%%%%%%%%%%%%%%%%%%%%%%%%%%%%%%%%%%%%%%%%%%%%%%%%%%%%%%%%%%%%

\begin{exercise}[1]
What are the two defining conditions for a symplectic form?
\end{exercise}

\begin{exercise}[1]
Write the standard symplectic form on

\[
\R^2.
\]
\end{exercise}

\begin{exercise}[1]
Write the standard symplectic form on

\[
\R^4.
\]
\end{exercise}

\begin{exercise}[1]
Write the standard symplectic form on

\[
\R^6.
\]
\end{exercise}

\begin{exercise}[1]
Explain why a symplectic manifold cannot have dimension \(5\).
\end{exercise}

\begin{exercise}[1]
Identify the position and momentum variables in

\[
(q_1,q_2,p_1,p_2).
\]
\end{exercise}

\begin{exercise}[2]
Show that

\[
dq\wedge dp
\]

is closed.
\end{exercise}

\begin{exercise}[2]
Show that

\[
dq\wedge dp
\]

is nondegenerate on \(\R^2\).
\end{exercise}

\begin{exercise}[2]
Compute

\[
\omega(u,v)
\]

for

\[
u=(1,2),
\qquad
v=(4,-1),
\]

using

\[
\omega=dq\wedge dp.
\]
\end{exercise}

\begin{exercise}[2]
Verify that the standard symplectic matrix

\[
J=
\begin{pmatrix}
0&1\\
-1&0
\end{pmatrix}
\]

satisfies

\[
J^2=-I.
\]
\end{exercise}

\begin{exercise}[2]
Explain the meaning of

\[
A^TJA=J.
\]
\end{exercise}

\begin{exercise}[2]
What does a symplectomorphism preserve?
\end{exercise}

\begin{exercise}[2]
State Darboux's theorem.
\end{exercise}

\begin{exercise}[2]
Explain why Darboux's theorem suggests that symplectic geometry has no local
curvature invariant associated solely with \(\omega\).
\end{exercise}

\begin{exercise}[2]
For

\[
H(q,p)=\frac{p^2}{2m}+V(q),
\]

derive Hamilton's equations.
\end{exercise}

\begin{exercise}[2]
For

\[
H(q,p)=\frac12p^2+\frac12q^2,
\]

find

\[
\dot q
\]

and

\[
\dot p.
\]
\end{exercise}

\begin{exercise}[2]
Compute

\[
\{q,p\}.
\]
\end{exercise}

\begin{exercise}[2]
Compute

\[
\{p,q\}.
\]
\end{exercise}

\begin{exercise}[2]
Compute

\[
\{q^3,p^2\}.
\]
\end{exercise}

\begin{exercise}[2]
Show that

\[
\{f,f\}=0.
\]
\end{exercise}

\begin{exercise}[2]
If

\[
\{f,H\}=0,
\]

what does this imply about \(f\) along the Hamiltonian flow when \(f\) has no
explicit time dependence?
\end{exercise}

\begin{exercise}[2]
If

\[
\dim Q=4,
\]

find

\[
\dim T^*Q.
\]
\end{exercise}

\begin{exercise}[2]
Write the canonical one-form on

\[
T^*\R^2.
\]
\end{exercise}

\begin{exercise}[2]
Use

\[
\theta=p\,dq
\]

to compute

\[
-d\theta.
\]
\end{exercise}

\begin{exercise}[2]
Define an isotropic subspace.
\end{exercise}

\begin{exercise}[2]
Define a Lagrangian subspace.
\end{exercise}

\begin{exercise}[2]
What is the dimension of a Lagrangian subspace of an
\(8\)-dimensional symplectic vector space?
\end{exercise}

\begin{exercise}[3]
Prove that a Lagrangian subspace is isotropic.
\end{exercise}

\begin{exercise}[3]
Explain why the \(q\)-coordinate subspace of \(\R^{2n}\) is Lagrangian.
\end{exercise}

\begin{exercise}[3]
Explain why the \(p\)-coordinate subspace is also Lagrangian.
\end{exercise}

\begin{exercise}[3]
Show that the graph of \(df\) is Lagrangian in \(T^*Q\).
\end{exercise}

\begin{exercise}[3]
Show that

\[
\omega^n
\]

is a top-degree form on a \(2n\)-dimensional symplectic manifold.
\end{exercise}

\begin{exercise}[3]
Explain why a symplectic manifold is orientable.
\end{exercise}

\begin{exercise}[3]
Explain why Hamiltonian flows preserve the Hamiltonian.
\end{exercise}

\begin{exercise}[3]
Explain the geometric meaning of Liouville's theorem.
\end{exercise}

\begin{exercise}[3]
Why is volume preservation weaker than symplectic preservation in dimensions
greater than two?
\end{exercise}

\begin{exercise}[3]
Explain the significance of Gromov's non-squeezing theorem.
\end{exercise}

\begin{exercise}[3]
What is a moment map intended to encode?
\end{exercise}

\begin{exercise}[3]
Explain the broad idea of symplectic reduction.
\end{exercise}

\begin{exercise}[3]
Explain how translation symmetry is related to momentum conservation.
\end{exercise}

\begin{exercise}[3]
Suppose \(J\) is compatible with \(\omega\). Explain why

\[
g(X,Y)=\omega(X,JY)
\]

behaves like a metric.
\end{exercise}

\begin{exercise}[3]
Explain why every K\"ahler manifold is symplectic.
\end{exercise}

\begin{exercise}[3]
Explain why the converse need not hold.
\end{exercise}

\begin{exercise}[3]
Prove that an exact form is closed.
\end{exercise}

\begin{exercise}[3]
Use Stokes' theorem to explain why a compact symplectic manifold without
boundary cannot have an exact symplectic form.
\end{exercise}

\begin{exercise}[4]
Research the symplectic structure of the two-sphere \(S^2\).
\end{exercise}

\begin{exercise}[4]
Research the symplectic structure on complex projective space

\[
\mathbb{P}^n_{\C}.
\]
\end{exercise}

\begin{exercise}[4]
Investigate the geometric meaning of a symplectic capacity.
\end{exercise}

\begin{exercise}[4]
Research Gromov's non-squeezing theorem and explain why ordinary volume
cannot detect the obstruction.
\end{exercise}

\begin{exercise}[4]
Investigate one example of symplectic reduction.
\end{exercise}

\begin{exercise}[4]
Research the moment map for rotational symmetry in classical mechanics.
\end{exercise}

\begin{exercise}[4]
Investigate the relationship between Lagrangian submanifolds and generating
functions.
\end{exercise}

\begin{exercise}[4]
Research the basic idea of Floer homology.
\end{exercise}

\begin{exercise}[4]
Investigate one application of symplectic geometry to celestial mechanics.
\end{exercise}

\begin{exercise}[4]
Research how symplectic integrators differ from ordinary numerical
integration methods.
\end{exercise}

%%%%%%%%%%%%%%%%%%%%%%%%%%%%%%%%%%%%%%%%%%%%%%%%%%%%%%%%%%%%
\subsection{Conceptual Summary}
\label{subsec:symplectic-summary}
%%%%%%%%%%%%%%%%%%%%%%%%%%%%%%%%%%%%%%%%%%%%%%%%%%%%%%%%%%%%

A symplectic manifold is a pair

\[
\boxed{
(M,\omega)
}
\]

where

\[
\omega
\]

is a closed, nondegenerate differential \(2\)-form.

Thus

\[
\boxed{
d\omega=0.
}
\]

The standard symplectic space is

\[
\boxed{
\left(
\R^{2n},
\sum_{j=1}^{n}dq_j\wedge dp_j
\right).
}
\]

Symplectic manifolds necessarily have even dimension.

Darboux's theorem states that every symplectic manifold is locally equivalent
to the standard model.

Therefore symplectic geometry has no local invariants associated solely with
the symplectic form analogous to Riemannian curvature.

A Hamiltonian

\[
H:M\to\R
\]

determines a Hamiltonian vector field \(X_H\).

With the convention used in this section,

\[
\boxed{
\iota_{X_H}\omega=dH.
}
\]

In canonical coordinates this produces Hamilton's equations

\[
\boxed{
\dot q_j
=
\frac{\partial H}{\partial p_j},
\qquad
\dot p_j
=
-
\frac{\partial H}{\partial q_j}.
}
\]

Hamiltonian flows preserve

\[
\boxed{
\omega,
}
\]

the Hamiltonian,

\[
\boxed{
H,
}
\]

and symplectic volume,

\[
\boxed{
\frac{\omega^n}{n!}.
}
\]

The Poisson bracket is

\[
\boxed{
\{f,g\}
=
\sum_j
\left(
\frac{\partial f}{\partial q_j}
\frac{\partial g}{\partial p_j}
-
\frac{\partial f}{\partial p_j}
\frac{\partial g}{\partial q_j}
\right).
}
\]

Conserved quantities satisfy

\[
\boxed{
\{f,H\}=0.
}
\]

Every cotangent bundle

\[
T^*Q
\]

has a canonical symplectic structure.

Lagrangian submanifolds have half the dimension of the ambient symplectic
manifold and satisfy

\[
\boxed{
\omega|_L=0.
}
\]

Moment maps connect symmetry with conserved quantities, while symplectic
reduction uses symmetry to construct smaller symplectic spaces.

Gromov's non-squeezing theorem demonstrates that symplectic geometry has
rigidity beyond ordinary volume preservation.

Finally, K\"ahler geometry provides a fundamental bridge among

\[
\boxed{
\text{symplectic},
\qquad
\text{complex},
\qquad
\text{Riemannian geometry}.
}
\]

%%%%%%%%%%%%%%%%%%%%%%%%%%%%%%%%%%%%%%%%%%%%%%%%%%%%%%%%%%%%
\subsection{Section Connections}
\label{subsec:symplectic-connections}
%%%%%%%%%%%%%%%%%%%%%%%%%%%%%%%%%%%%%%%%%%%%%%%%%%%%%%%%%%%%

\chapterconnection{
Symplectic Geometry builds directly on Differential Geometry through
differential forms and smooth manifolds and on Complex Geometry through
compatible almost complex and K\"ahler structures. Its canonical
position-momentum coordinates provide the geometric foundation of
Hamiltonian mechanics and connect naturally with Differential Equations and
Dynamical Systems. Poisson brackets connect symplectic geometry with Lie
algebras, while moment maps and symplectic reduction connect the subject with
group actions and symmetry. Lagrangian submanifolds and symplectic
invariants lead toward modern topology, Floer theory, and mirror symmetry.
The geometric behavior of convex subsets of symplectic vector spaces also
provides a natural transition to the next section, Convex Geometry.
}

%%%%%%%%%%%%%%%%%%%%%%%%%%%%%%%%%%%%%%%%%%%%%%%%%%%%%%%%%%%%
% SECTION SOURCE TRACKING
%%%%%%%%%%%%%%%%%%%%%%%%%%%%%%%%%%%%%%%%%%%%%%%%%%%%%%%%%%%%

%------------------------------------------------------------
% 3.10 Symplectic Geometry
%------------------------------------------------------------
%
% Source basis:
%   - Standard introductory symplectic geometry,
%     differential geometry, and Hamiltonian mechanics.
%
% Used for:
%   - symplectic forms and manifolds;
%   - closedness and nondegeneracy;
%   - standard symplectic structures;
%   - canonical position-momentum coordinates;
%   - symplectic vector spaces and bases;
%   - symplectic matrices and transformations;
%   - symplectomorphisms;
%   - Darboux's theorem;
%   - phase space;
%   - Hamiltonian functions and vector fields;
%   - Hamilton's equations;
%   - Hamiltonian flows;
%   - Poisson brackets;
%   - conservation laws;
%   - cotangent bundles;
%   - canonical one-forms;
%   - isotropic, coisotropic, and Lagrangian geometry;
%   - symplectic area and volume;
%   - Liouville's theorem;
%   - non-squeezing;
%   - moment maps;
%   - symplectic reduction;
%   - compatible almost complex structures;
%   - Kahler geometry connections;
%   - exact symplectic forms;
%   - contact geometry preview.
%
% Notes:
%   - Differential forms, tangent bundles, Lie derivatives,
%     and smooth manifolds are developed canonically in
%     Differential Geometry.
%   - Complex and Kahler structures are developed in
%     Complex Geometry.
%   - Hamiltonian differential equations are revisited in
%     Differential Equations and Dynamical Systems.
%   - Lie groups, Lie algebras, and group actions are
%     developed canonically in Algebra and Geometry.
%   - Floer theory and advanced symplectic topology are
%     deferred to later advanced sections.
%   - Sign conventions for Hamiltonian vector fields vary
%     across the literature; this section uses
%     i_{X_H} omega = dH with omega = sum dq_j wedge dp_j.
%
%%%%%%%%%%%%%%%%%%%%%%%%%%%%%%%%%%%%%%%%%%%%%%%%%%%%%%%%%%%%